\documentclass[../main.tex]{subfiles}
\begin{document}
%==========================================TIÊU ĐỀ TẬP========================================
\begin{table}[h]
  \begin{adjustwidth}{-1cm}{}
    \begin{tabular}{>{\centering\arraybackslash}p{6cm}|>{\raggedright\arraybackslash}p{12.5cm}}
      \multirow{5}{6cm}
      % Cột bên trái: ÉPISODE và số tập
      {\\[-40pt]
        \flaregothic\fontsize{20pt}{20pt}\selectfont \centering
        \color{eptype}ÉPISODE\color{black}\\
        \burbank\fontsize{72pt}{72pt}\selectfont
        \color{epnum}45\color{black}\\[6pt]
        \cafeteria\fontsize{20pt}{20pt}\selectfont\color{epnum}(4-3)\color{black}
        \\[18pt]
        % \flaregothic\fontsize{14pt}{14pt}\selectfont
        % \color{epattr}BIENVENUE, \uline{\color{Banpire} Mel} !
        \flaregothic\fontsize{18pt}{18pt}\selectfont
        \color{epattr}ÉPISODE PERDU
        \color{black}
      }
       &
      % Dòng thứ nhất: tên tập tiếng Nhật
      {\fotskip\fontsize{18pt}{18pt}\selectfont \ruby{嫌}{いや}がらせを\ruby{受}{う}けています
      }       \\[6pt]
      % Dòng thứ hai: tên tập tiếng Anh
       & {\cafeteria\fontsize{18pt}{18pt}\selectfont This is the Curry Land}          \\[4pt]
      % Dòng thứ ba: tên tập tiếng Việt
       & {\vnmsans\fontsize{18pt}{18pt}\selectfont Buổi ăn cà ri khó nhọc của Roboco} \\[8pt]
      % Dòng thứ tư: Nhân vật xuất hiện
       & & {\caviar{\fontsize{14pt}{14pt}\selectfont Nhân vật: \emp{kuon1} \emp{roboco2} \emp{fubuki} \emp{matsuri} \emp{pekora}}}
      \\[8pt]
      % Dòng cuối cùng: Thời gian diễn ra của tập (thường sẽ là ngày phát hành)
       & {\continuum\fontsize{16pt}{16pt}\selectfont Ngày phát hành: 15/03/2020}      \\[18pt]
    \end{tabular}
  \end{adjustwidth}
\end{table}
%=========================================TRUYỆN CHÍNH========================================
\color{black}

\txtbx[2]{{\color{vol} \uline{Tóm tắt:}} Một cơn thèm cà ri bất chợt của Roboco đã kéo cô nàng vào một chuỗi những sự kiện dở khóc dở cười. Sau khi mua nhầm hộp cà ri ``Sina Dam'' rởm chứa hai ông bà cụ tí hon biết đọc tiên tri, cô nàng Người máy quyết định tự tay vào bếp dưới sự động viên của Kuon.}

Một ngày trung tuần tháng Ba rực rỡ, hơi thở dịu dàng của mùa xuân đang dần tràn ngập khắp từng ngóc ngách của thành phố Điện tử sầm uất. Tiết trời mỗi ngày một thêm phần ấm áp, dễ chịu, cái se se lạnh buốt giá của những ngày đông hỗn loạn dường như đã không còn đậm nét trong văn phòng mới nữa.

Thế nhưng, dù là mùa đông hay mùa xuân, đối với một cô nàng Người máy như Roboco, sự thay đổi của thời tiết hoàn toàn không phải là điều quá quan trọng. Điều duy nhất đang chiếm trọn tâm trí và khiến cô nàng đứng ngồi không yên lúc này là\dots

Cô nàng khẽ áp hai tay lên má, đôi mắt sáng rực lên mơ màng: ``Tự dưng mình lại thèm ăn một đĩa cà ri bốc khói nghi ngút quá đi mất\dots\ Chắc phải tranh thủ chạy ra ngoài mua một hộp mới được\dots''

Nàng Người máy bỗng dưng thèm ăn cà ri một cách mãnh liệt. Ánh mắt cô nàng lấp lánh tia hy vọng khi tưởng tượng đến mùi hương thơm nức mũi của bột gia vị hòa quyện cùng vị cay nồng, ấm áp đặc trưng của món ăn.

Khác hẳn với phong cách ăn mặc tối giản ngày thường, hôm nay Roboco khoác thêm một chiếc áo hoodie trùm đầu màu hồng mang họa tiết rằn ri quân đội vô cùng thời thượng, bao trọn nửa thân trên của cô nàng. Thế nhưng, phần thân dưới của cô thì vẫn y như nguyên trạng: hoàn toàn không hề mặc thêm bất kỳ chiếc quần nào cả. Mái tóc của cô nàng hôm nay đã được ``nuôi'' dài ra (thực chất ra là cô nàng vừa mới lắp đặt thêm bộ tóc giả mới cao cấp), có thể tự do biến tấu, buộc thành nhiều kiểu dáng khác nhau. Tuy nhiên hôm nay, cô quyết định chọn kiểu buộc tóc đuôi ngựa đơn giản, năng động kết hợp đeo một chiếc kính gọng tròn không tròng đầy trí thức.

\model{27}{l}{6cm}{roboco2}

Dẫu biết rõ cô nàng là một rô-bốt sinh học nên cơ thể hoàn toàn không có cảm giác lạnh là gì, nhưng có lẽ cô nàng cũng muốn được diện đồ thời trang xuân một chút cho bằng chị bằng em. Hoặc cũng có thể là ai đó trong ban quản lý đã khuyên cô nên mặc thêm chiếc áo ngoài để trông bớt hở hang hơn, dẫu rằng chiếc áo mỏng manh ấy cũng chẳng che chắn thêm được bao nhiêu phần nhạy cảm cả.

Kuon đang ngồi kiểm tra đống sổ sách gần đó nghe thấy lời thì thầm đầy thèm thuồng của cô liền khẽ ngẩng đầu lên khỏi đống giấy tờ, nhiệt tình đề xuất: ``Này Roboco-chan, anh thấy trong bếp công ty mình vẫn còn lưu kho đầy đủ các nguyên liệu cơ bản để nấu cà ri đấy chứ. Sao em không tự tay vào bếp làm một nồi tại đây luôn cho tiết kiệm, lại đảm bảo vệ sinh nữa?''

Roboco nghe thấy thế liền tỏ ra hứng khởi và thích thú hơn hẳn, hai tai khẽ động đậy: ``Thật thế hả anh? Nếu trong bếp có sẵn đồ thì tốt quá, đỡ mất công em phải lội bộ ra ngoài bách hóa mua rồi! Mà\dots\ đống nguyên liệu đó anh cất ở ngăn tủ nào thế?''

Thế nhưng, ngay trước khi Kuon kịp mở lời chỉ chỗ cho cô, cánh cửa gỗ phòng họp bỗng bật mở kèm theo một giọng nói vô cùng quen thuộc vang lên. Fubuki chạy vội vào trong với dáng vẻ hớt hải, mặt mũi tái mét như thể vừa mới gặp phải một chuyện vô cùng hệ trọng cần nhờ vả cô nàng Người máy kia gấp. 

Cô nàng vừa thở dốc vừa lớn tiếng gọi cứu viện: ``A! Roboco-san ơi! Cứu em với, cứu em xử lý mấy củ hành tây quái đản kia với ạ!''

Nàng Người máy chớp chớp mắt đầy bối rối, khẽ nghiêng đầu hỏi lại với vẻ mặt ngơ ngác: ``Hành tây sao em? Mấy củ hành tây thì có chuyện gì làm sao cơ?''

Nàng Cáo trắng vẫn thở hổn hển không ra hơi, vội vàng hoa tay múa chân giải thích: ``Thì là đống hành tây lần trước em lỡ tay để quên ngoài nắng mấy ngày liền ấy mà! Bây giờ tự dưng trên thân tụi nó mọc đầy một lớp lông trắng muốt nhìn dị hợm lắm luôn! Tụi mình phải nhanh chóng dọn dẹp và xử lý gấp đống đó đi thôi, chứ không là tụi nó sắp sửa 'xâm chiếm' toàn bộ căn bếp của văn phòng tụi mình đến nơi rồi đấy!''

Để làm tăng thêm mức độ nghiêm trọng và thuyết phục cho câu chuyện của mình, cô nàng Cáo trắng còn không quên thêm thắt chút gia vị ``màu mè'' vào lời kể. Cô vừa nói vừa cố tình căng hai cánh tay lên, tạo dáng gồng cơ bắp như một vận động viên thể hình lực điền thực thụ khiến cho Roboco lại càng thêm phần ngơ ngác và kinh ngạc trước màn biểu diễn đầy kịch tính này.

\img{4-3a}

Chứng kiến màn gồng cơ bắp có phần vô tri của Fubuki nhưng lại không nhận được bất kỳ phản ứng kinh ngạc nào từ hai người nghe, đôi tai cáo mềm mại của cô nàng liền rũ cụp xuống đầy chán nản, dáng đứng cũng xụ xuống hệt như một quả bóng bị xì hơi. 

Kuon đứng bên cạnh khẽ thở dài đầy bất lực, cất giọng châm chọc: ``Nói thật là anh nãy giờ vẫn chưa nhìn thấy cái cơ bắp lực điền nào của em đâu đấy nhé Fubu-chan\dots\ Mà có chuyện gì nhờ vả thì cứ mở miệng nói thẳng ra đi, lại còn bày đặt diễn sâu kể lể kịch tính như phim hành động viễn tưởng thế kia nữa chứ!''

Roboco cũng khẽ bật cười trừ trước điệu bộ đáng yêu của người em, nhẹ nhàng hỏi: ``Thế rốt cuộc là có chuyện gì cần chị giúp thế?''

Không còn cách nào khác để diễn kịch nữa, Fubuki đành chỉ tay về phía chiếc thùng gỗ lớn đặt ở góc cửa, ỉu xìu nói: ``Dạ, hai người giúp em khênh cái thùng hành tây này chuyển qua bên phòng kho với ạ! Chứ cứ để nó nằm xó ngoài hành lang thế này thêm mấy ngày nữa thì chắc tụi nó mọc rễ bám chặt xuống sàn nhà luôn quá!''

Nàng Người máy nhìn chiếc thùng hành tây đang đặt ở góc cửa, khẽ nhíu mày thắc mắc: ``Chị nhìn cái thùng kia trông cũng nhỏ nhắn, nhẹ hều mà ta. Sao tự dưng một cô cáo khỏe mạnh như em lại không tự khênh đi mà phải nhờ đến chị làm gì nhỉ?''

Roboco còn đang mải mê phân tích logic thì cậu Tổng quản lý Kuon đã nhanh nhảu xen vào để động viên cô nàng: ``Thôi nào Roboco-chan, coi như đây là một bài tập thể dục khởi động nhẹ nhàng, giãn gân cốt trước khi bắt tay vào công việc nấu nồi cà ri thơm ngon đi.''

Nghe cậu Tổng nói cũng có lý, nàng Người máy chỉ biết khẽ thở dài một hơi đầy bất lực. Cô quay sang nhìn gương mặt đang tràn đầy sự mong chờ và cầu khẩn của Fubuki gật đầu đồng ý: ``Thôi được rồi, được rồi\dots\ Để chị giúp em một tay vậy.''

Thế là, hai người họ bắt đầu khom lưng, hì hục cùng nhau phối hợp nâng chiếc thùng hành tây khổng lồ lên để di chuyển ra phòng kho. Thế nhưng, ngay khi vừa mới nhấc chiếc thùng lên khỏi mặt đất, cả Roboco lẫn Kuon đều lập tức nhận ra trọng lượng của chiếc thùng này nặng hơn rất nhiều so với những gì họ hằng tưởng tượng. 

Kuon cố gắng bậm môi giữ thăng bằng, lẩm bẩm đầy nghi hoặc: ``Này Fubuki ơi, cái đống này không phải là hành tây bình thường đúng không? Sao nó lại nặng đến mức vô lý thế này cơ chứ?''

Nàng Cáo trắng đứng bên cạnh chỉ biết cười khúc khích đầy ranh mãnh, nhất quyết không chịu mở miệng trả lời thẳng vào câu hỏi. Nàng Người máy thấy thái độ mập mờ ấy liền tỏ ra nghi ngờ hỏi: ``Không lẽ nào\dots\ em lại lén lút nhét thêm cái thứ gì nặng nề vào trong đó nữa rồi đúng không?''

Fubuki gãi gãi đầu cười khì khì đầy vô tội giải thích: ``À thì\dots\ thực ra em cũng có đổ thêm một chút đất xốp vào bên trong thùng hành tây ấy mà\dots\ mục đích chính là để giữ cho hành tây luôn có độ ẩm và tươi lâu hơn thôi chứ không có ý gì khác đâu nha!''

Cả Roboco lẫn Kuon nghe xong lời thú nhận xanh rờn ấy liền đứng hình mất vài giây toàn tập. Hai người đưa mắt nhìn nhau đầy bất lực, rồi lại nhìn chằm chằm vào cái thùng hành tây bám đầy đất xốp siêu nặng như thể vừa mới phát hiện ra một bí mật kinh thiên động địa nào đó vậy.

\ct

Dẫu trong lòng đầy bất bình trước trò nghịch ngợm của nàng Cáo, cô nàng Roboco và Kuon vẫn phải cùng nhau gồng sức nhấc chiếc thùng hành tây nặng trịch lên. Hai người vừa khệ nệ bước đi từng bước ngắn, vừa phải cố gắng giữ thăng bằng tuyệt đối vì chiếc thùng khá to lớn và cồng kềnh, chắn hết cả tầm nhìn phía trước. Fubuki đứng gần đó liền nhanh nhảu ngồi xổm xuống, đưa mắt theo dõi với vẻ mặt vô cùng lo lắng và hồi hộp. Dù bản thân hoàn toàn không phải là người trực tiếp tham gia lao động nặng nhọc, cô nàng vẫn vô cùng nhiệt tình cất giọng cổ vũ, hò hét inh ỏi để khích lệ tinh thần đồng đội: ``Cố lên hai người ơi! Một chút xíu nữa thôi là sắp đến nơi rồi!''

\img{4-3b}

Kuon tức mình liếc xéo Fubuki một cái rõ sắc lẻm nhưng đang bận giữ chiếc thùng nặng nên chẳng buồn mở miệng đáp lại. Trong lúc cả hai người đang cùng nhau dùng lực để nâng chiếc thùng vượt qua một bậc cửa cao, họ liền đồng thanh hô vang khẩu hiệu để lấy đà: ``Một, hai, ba\dots\ A-lê-hấp!''

Roboco khẽ ngẩng gương mặt lấm tấm mồ hôi lên, tự tin tuyên bố: ``Mà này Kuon-senpai ơi, thực ra cái thùng này một mình em dư sức tự khênh đi được rồi mà. Anh không cần phải tốn sức vào giúp đỡ em đâu.''

Cậu con trai nghe thấy thế liền khẽ nhếch môi nở một nụ cười đầy châm chọc. Cậu vẫn giữ chặt lấy tay cầm của chiếc thùng và đáp lại bằng giọng điệu trêu đọc quen thuộc: ``Kệ anh chứ, anh thích vào giúp đỡ em thì anh cứ giúp thôi, đồ Robo-PON ngốc nghếch ạ. Ngộ nhỡ tay em lại rụng xuống như hồi năm ngoái\footnote{Xem lại {\flaregothic ÉPISODE 5}, quyển số Một.} thì sao.''

Vừa nghe thấy danh xưng chọc ghẹo của mình từ miệng cậu Tổng, cô nàng Người máy lập tức nổi quạu. Cô quay sang trừng mắt nhìn cậu bằng ánh mắt đầy bực tức, hét lên phản đối dữ dội: ``Sao lúc nào anh cũng tỏ thái độ không tin tưởng vào sức mạnh của em thế nhở!! Em đã bảo là em hoàn toàn đủ mạnh mẽ để có thể tự mình giải quyết đống này mà lị!!''

Kuon nhún vai cười trừ đáp lại: ``Sau mấy cái chiến tích huy hoàng phá tung tanh bành cả văn phòng công ty của em trước đây thì anh nghĩ mình chẳng còn đủ dũng khí để tin tưởng vào lời hứa của em nữa đâu.'' Câu trả lời vô cùng thực tế và phũ phàng của cậu Tổng lại càng khiến cho cô nàng Người máy thêm phần tức tối, phồng má giận dỗi.

Fubuki nhìn thấy bầu không khí bắt đầu có dấu hiệu căng thẳng dở khóc dở cười liền vội vàng xua xua tay, cất giọng cố gắng đứng ra giảng hòa: ``Thôi nào hai người ơi, đừng cãi cọ nhau trong lúc đang khênh đồ nặng thế chứ! Cố gắng nhích thêm một chút xíu nữa thôi, sắp đặt được thùng vào phòng kho rồi kìa!''

Cuối cùng, sau bao nỗ lực vất vả gồng gánh cơ bắp, chiếc thùng hành tây siêu nặng cũng đã được đặt gọn gàng, ngăn nắp vào một góc khuất bên trong phòng kho. Nàng Cáo trắng Fubuki không giấu nổi niềm vui sướng tột độ, cô nàng rối rít vỗ tay rộn ràng và cất giọng cảm ơn đầy chân thành: ``Ôi, cảm ơn hai người nhiều nhiều lắm nha! Nếu hôm nay không có sự giúp đỡ đắc lực của hai người thì chắc em chỉ biết khóc ròng chịu chết bên cạnh cái thùng gỗ siêu nặng này mất thôi!''

Roboco đưa tay phủi phủi bụi bẩn trên quần áo, khẽ thở phào nhẹ nhõm đầy sướng khoái: ``Phù, cuối cùng cũng xong việc rồi đấy nhỉ, tốn mất của tụi mình một khoảng thời gian dài rồi đấy Kuon-senpai.''

Kuon trái lại vẫn giữ nguyên vẻ mặt vô cùng điềm nhiên và thản nhiên như không có chuyện gì, khẽ nhún vai đáp: ``Làm gì mà mất nhiều thời gian đến thế chứ em. Cơ mà nếu để một mình em tự bê cái thùng nặng này thì có khi giữa đường đã rụng mất cả hai cánh tay máy ra rồi ấy chứ.''

Roboco lườm cậu Tổng một cái rõ sắc lẻm hệt như muốn phóng ra dao găm, nhưng do cơ thể đã bắt đầu thấm mệt nên cô nàng cũng chẳng buồn mở miệng cãi cọ lại làm gì cho tốn sức. Thay vào đó, cô nàng khẽ đưa tay vỗ nhẹ nhẹ vào chiếc bụng nhỏ của mình, cất giọng than thở đầy đói lã: ``Mà thôi, dẹp chuyện đó đi, giờ em đói bụng rã rời cả người rồi đây này. Bây giờ em phải mau chóng chạy đi mua một hộp cà ri ăn để nạp lại năng lượng gấp mới được--''

Thế nhưng, nàng Người máy còn chưa kịp quay người bước ra khỏi phòng kho để đi mua đồ ăn thì bỗng nhiên cảm nhận được một cánh tay thình lình đặt mạnh lên vai mình từ phía sau. Cô giật mình hoảng hốt, vội vàng quay phắt người lại nhìn thì đập vào mắt là gương mặt đang đầy vẻ mếu máo và cầu cứu của Matsuri. Giọng nói của cô nàng Lễ hội vang lên đầy khẩn thiết, dồn dập: ``Giúp em với, Roboco-senpai ơi! Làm ơn cứu mạng em với!''

\img{4-3c}

Kuon đứng tựa người bên cạnh chứng kiến cảnh tượng quen thuộc này liền khẽ lắc đầu cười trừ, thản nhiên buông một câu trêu chọc đầy phũ phàng: ``Ái chà, xem ra cái ước mơ được ăn đĩa cà ri nóng hổi của em lại phải dời lại vô thời hạn rồi nhé Roboco-chan.''

Roboco nghe xong lời chọc ghẹo của cậu Tổng cũng chỉ biết nở một nụ cười khổ đầy bất lực, khẽ gật đầu đồng tình: ``Chắc là số phận bắt em phải nhịn đói rồi anh ạ.'' Cô quay sang nhìn Matsuri, khẽ chống hai tay ngang hông hỏi thăm: ``Có chuyện gì thế Matsuri-chan? Sao trông em cuống cuồng hết cả lên thế kia?''

Nàng Cổ động liền lập tức trình bày hoàn cảnh khó khăn của mình. Cô nàng đưa hai tay ra, trưng ra một vật thể kỳ lạ trông vô cùng giống một củ khoai tây bình thường đang cầm trên tay: ``Là bức tượng khoai tây kỷ niệm vô cùng quý giá của em tự dưng hôm nay lại dở chứng không chịu hoạt động nữa rồi đây này! Chị Roboco-senpai nổi tiếng là thiên tài công nghệ máy móc, chị có giúp em kiểm tra và sửa chữa nó được không ạ?''

\img{4-3d}

Kuon nheo nheo mắt nhìn kỹ vào vật thể lạ đang nằm gọn trong lòng bàn tay của Matsuri. Dưới góc nhìn thực tế của cậu, cái thứ này trông hoàn toàn không có bất kỳ điểm nào khác biệt so với một củ khoai tây bình thường cả. Cậu khẽ nhún vai, đưa ra lời nhận xét đầy thẳng thắn: ``Anh nhìn đi nhìn lại thì thấy cái này rõ ràng là một củ khoai tây thật mà em. Làm gì có cái chi tiết máy móc hay nút bấm nào ở trên đó đâu chứ?''

Nàng Người máy cũng tò mò cúi thấp người xuống để xem xét củ khoai một cách kỹ lưỡng hơn. Cô nàng khẽ lắc đầu tỏ vẻ ngạc nhiên đầy hoang mang: ``Chị thực sự cũng không thể hiểu nổi cái bức tượng củ khoai này hoạt động theo cơ chế vật lý nào nữa luôn đấy\dots\ Mà nhìn qua thì chị có thấy nó bị nứt vỡ hay hỏng hóc ở chỗ nào đâu em?''

Matsuri lắc đầu liên tục như một chiếc máy, ra sức giải thích với gương mặt bắt đầu mất đi sự bình tĩnh: ``Không phải đâu mà! Bình thường chiếc tượng khoai tây này hay phát ra tiếng reo chuông báo thức vui tai lắm luôn á! Chắc chắn là hệ thống máy móc phức tạp ẩn giấu bên trong tụi nó bị hư hỏng ở đâu đó rồi!''

Cậu Tổng quản lý tò mò hỏi tiếp: ``Thế nó bắt đầu bị tình trạng 'tịt ngòi' không thèm kêu thế này từ bao giờ rồi em?''

``Dạ, mới bị từ sáng sớm ngày hôm nay thôi ạ!'' Matsuri buồn rầu đáp, đồng thời tiếp tục hướng đôi mắt long lanh đầy vẻ khẩn cầu về phía hai người.

Roboco chăm chú nhìn Matsuri, rồi lại quay sang nhìn củ khoai tây kỳ lạ kia, gương mặt cô nàng đăm chiêu, nhíu mày suy nghĩ tìm giải pháp xử lý. Cậu con trai đứng bên cạnh, hai tay khoanh trước ngực, vừa quan sát biểu cảm của hai cô gái vừa nhướng nhướng mày trêu chọc: ``Nói thật là anh vẫn thấy nó giống một củ khoai tây để nấu súp hơn là một bức tượng công nghệ đấy\dots\ Nhưng mà nhìn thái độ của em thì chắc là món đồ này mang một ý nghĩa kỷ niệm vô cùng quan trọng đối với em lắm nhỉ?''

Nàng Cổ động gật đầu lắt léo đầy khẳng định: ``Vâng ạ! Nó là món quà kỷ niệm vô cùng quý giá và độc nhất vô nhị của em đấy ạ!''

Kuon nhìn lại củ khoai tây đó thêm một lần nữa, gương mặt vẫn không giấu nổi vẻ hoài nghi sâu sắc: ``Anh thì vẫn không thể nào hiểu nổi mức độ quý giá của một củ khoai tây nằm ở chỗ nào, nhưng thôi, nếu em đã khẳng định như thế thì\dots''

Nàng Người máy khẽ nhắm nghiền đôi mắt lại, thở dài một hơi dài đầy cam chịu. Lúc này trong đầu cô nàng đang diễn ra một trận chiến tư tưởng vô cùng quyết liệt giữa một bên là cơn đói bụng đang cồn cào réo rắt đòi ăn cà ri, và một bên là mong muốn được giúp đỡ cô em đồng nghiệp Matsuri đang gặp khó khăn. Cô thầm nghĩ trong bụng: ``\textit{Trời đất ơi, sắp đến giờ ăn trưa rồi mà bụng mình thì đói run cả người lên rồi. Thế nhưng Matsuri-chan đang vô cùng cần sự trợ giúp công nghệ của mình để sửa món đồ kỷ niệm của em ấy, mình không thể nhẫn tâm từ chối được\dots}''

Cô khẽ liếc mắt nhìn sang phía cậu Tổng quản lý đang đứng thảnh thơi bên cạnh, thầm suy tính: ``\textit{Hay là mình thử mở lời nhờ Kuon-senpai cùng vào giúp một tay nhỉ? Thế nhưng liệu một người chuyên ngành IT như anh ấy có biết sửa chữa mấy cái món đồ máy móc cơ khí kết hợp thủ công kỳ dị kiểu này không ta?}''

\img{4-3e}

Cô nàng đưa mắt liếc nhìn chiếc đồng hồ treo tường lớn ở hành lang, thấy kim giờ đã chỉ sang cột mốc một giờ trưa từ lúc nào. Cô thầm tặc lưỡi tự nhủ: ``\textit{Thôi kệ đi vậy, đằng nào thì hôm nay mình cũng đã bị lỡ bữa ăn trưa mất rồi. Thay vì cứ đứng đây phân vân thì thà cứ vui vẻ ra tay giúp đỡ em ấy một cúp cho xong chuyện vậy!}''

Cuối cùng, Roboco khẽ mở mắt ra, nở một nụ cười nhẹ nhàng và gật đầu đồng ý với Matsuri: ``Được rồi, chị quyết định sẽ giúp em sửa chữa bức tượng này. Thế nhưng nói trước là có vẻ hệ thống bên trong nó khá phức tạp nên sẽ phải tốn một khoảng thời gian kha khá đấy nhé.''

Kuon đứng bên cạnh cũng tiến lại gần, xắn tay áo lên cười đáp: ``Để anh cùng vào giúp em một tay luôn cho nhanh nhé Roboco-chan. Dù sao thì anh đây cũng có chút kiến thức và am hiểu về sơ đồ mạch điện tử đấy.'' Vừa dứt lời, bằng một cách thần kỳ nào đó, trên tay cậu Tổng đã xuất hiện sẵn một bộ dụng cụ sửa chữa mini chuyên nghiệp đầy đủ tua-vít và kìm hàn.

Nghe thấy cả hai người đồng ý giúp đỡ, Matsuri liền vui mừng khôn xiết, rối rít cúi đầu cảm ơn liên tục không ngừng, đôi mắt sáng rực lên đầy hy vọng rạng rỡ. Roboco cười nhẹ đầy chiều chuộng, bắt tay vào việc cầm lấy củ khoai tây để tiến hành xem xét kỹ lưỡng, trong khi Kuon đứng sát bên cạnh, lặng lẽ quan sát từng chuyển động của cô nàng Người máy.

\ct

Cả hai người họ cùng nhau ngồi ghé sát đầu lại, cẩn thận cúi xuống chăm chú xem xét cấu trúc của ``củ khoai tây'' kỳ lạ của Matsuri. Sau khi quan sát kỹ dưới kính lúp, hóa ra cái thứ trông giống hệt củ khoai tây thật này thực chất là một chiếc đồng hồ báo thức công nghệ cao được thiết kế lớp vỏ bọc vô cùng tinh vi để ngụy trang. Matsuri đứng sát bên cạnh, hồi hộp chỉ tay vào từng chi tiết nhỏ giải thích: ``Dạ, chỗ này hình như là cái nút bấm khởi động bị kẹt cứng không ấn xuống được nè\dots\ còn ở phía bên hông này nữa, em có cảm giác như hệ thống dây dẫn nối nguồn bên trong bị đứt rời ra rồi hay sao ấy ạ.''

Kuon nheo mắt quan sát, cẩn thận dùng nhíp để kiểm tra từng lỗi hỏng hóc mà chủ nhân của chiếc đồng hồ vừa mới chỉ ra. Roboco nhanh nhẹn cầm lấy một chiếc tua-vít mini từ hộp dụng cụ, chậm rãi và khéo léo vặn mở từng chiếc ốc vít nhỏ để bóc tách lớp vỏ ngoài giả khoai tây của chiếc đồng hồ báo thức ra. Đúng như dự đoán, ẩn giấu bên trong lớp vỏ là một hệ thống dây điện chằng chịt kết hợp với các bánh răng cơ khí hoạt động khá đơn giản nhưng đã bị bám đầy bụi bẩn.

Roboco khẽ gật gù nhận xét: ``Ồ, nhìn cấu trúc bên trong thế này thì có vẻ như hệ thống máy móc của nó cũng không quá phức tạp như chị hằng tưởng tượng đâu nhỉ.''

Kuon gật đầu đồng ý, tay khéo léo dùng mỏ hàn nhiệt để nối lại đoạn dây điện bị đứt: ``Ừ, anh cũng nghĩ là chỉ cần nối lại cái dây nguồn này là chiếc đồng hồ hoạt động bình thường trở lại được thôi. Mà này Mat-chan, cái món đồ chơi độc lạ này rốt cuộc là em mua ở đâu thế hả em?''

Nàng Cổ động ngượng ngùng gãi gãi đầu, đáp lại bằng chất giọng có chút hối thúc, ngượng ngùng: ``Dạ, cái này là em được người ta tặng cho trong một buổi lễ hội mùa hè năm ngoái đấy ạ. Hình như là phần thưởng giành được từ một trong những gian hàng trò chơi ném lon trúng thưởng ở hội chợ thì phải ạ.''

Roboco mỉm cười gật gù, đôi tay máy vẫn tiếp tục tỉ mỉ lắp ráp lại các bánh răng: ``Thảo nào, bảo sao chất lượng gia công bên trong của nó có vẻ không được bền bỉ cho lắm. Lần sau sau khi sửa xong thì nhớ giữ gìn cẩn thận hơn một chút nhé, Matsuri-chan.''

Cả hai người tiếp tục công việc sửa chữa trong bầu không khí im lặng, tập trung cao độ. Sau một khoảng thời gian loay hoay căn chỉnh, Roboco khéo léo lắp ráp lại phần vỏ ngoài giả khoai tây của chiếc đồng hồ báo thức rồi khẽ nhấn nút khởi động thử nghiệm. Ngay lập tức, từ bên trong ``củ khoai tây'' bỗng phát ra những tiếng chuông reo báo thức vô cùng vui tai, giòn giã, đồng thời toàn bộ củ khoai còn liên tục rung lắc bần bật đầy sinh động trên mặt bàn họp.

Matsuri nhảy cẫng lên đầy phấn khích, gương mặt rạng rỡ nụ cười: ``Yatta! Hoạt động lại bình thường rồi nè! Ôi, cảm ơn chị Roboco-senpai và anh Kuon nhiều nhiều lắm luôn á!''

Cô nàng Rô-bốt mỉm cười xoa nhẹ trán, cảm thấy hơi mệt sau khi tập trung cao độ: ``Thế này là đã hoàn toàn ổn thỏa theo đúng ý em rồi chứ hả cô bé?''

Nàng Lễ hội hớn hở ôm chặt lấy món đồ chơi yêu thích của mình vào lòng, gật đầu lia lịa: ``Dạ vâng ạ, hoạt động mượt mà và ổn định lắm luôn rồi ạ!''

Ngay sau khi hoàn thành xong xuôi công việc cứu hộ công nghệ, Roboco lập tức đứng phắt dậy, phủi phủi hai bàn tay máy của mình rồi hướng đôi mắt quyết tâm về phía cửa phòng họp: ``Được rồi! Nhiệm vụ cứu hộ đã hoàn thành xuất sắc, bây giờ thì chị phải lập tức chạy ra ngoài bách hóa tổng hợp để mua bằng được một hộp cà ri ăn cứu đói đây!''

\img{4-3f}

Cô nàng nói xong, không chờ hai người kia kịp phản ứng hay mở lời giữ lại, đã nhanh chóng xoay người mở toang cánh cửa phòng họp rồi co giò chạy biến mất hút dọc hành lang, còn không quên cẩn thận dùng lực đóng sầm cánh cửa lại phía sau phát ra tiếng động rõ lớn. Matsuri ngơ ngác chớp chớp mắt kinh ngạc, khẽ vẫy tay chào theo bóng lưng đã biến mất của đàn chị: ``Ể? Chị ấy chạy nhanh thế\dots\ Vậy, gặp chị sau nhé Roboco-senpai.''

Kuon khẽ thở dài một hơi đầy bất lực, vừa nhìn cánh cửa đóng chặt vừa lắc đầu cười khổ: ``Lạ lùng thật đấy chứ. Ban nãy anh đã nói rõ ràng là trong bếp công ty mình có sẵn đầy đủ nguyên liệu để nấu cà ri rồi cơ mà nhỉ? Tại sao con bé ấy vẫn cứ nhất quyết phải chạy ra ngoài bách hóa mua bằng được hộp cà ri ăn liền về làm gì cho mệt thế không biết?''

Nàng Lễ hội cười khúc khích đầy ranh mãnh, vẫn đang mân mê chiếc máy báo thức: ``Có thể chị ấy thích ra ngoài ăn hơn chăng?''

Cậu con trai chỉ biết lắc đầu thở dài thêm một lần nữa trước suy nghĩ ngây ngô đó, nhưng rồi khóe môi cũng khẽ cong lên thành một nụ cười vì sự tự do vô cùng đáng yêu của cô nàng Người máy: ``\textit{Haiz, chắc là con bé sợ ở lại đây để bị quấy rầy nữa đây mà.}''

Cậu liếc mắt nhìn sang chiếc đồng hồ treo tường, ánh mắt khẽ giật mình: ``\dots Quái, đến giờ này rồi à?''

\ct

Ngoài trời lúc này màn đêm u tối đã hoàn toàn buông xuống từ lâu, ánh đèn đường neon rực rỡ bắt đầu thắp sáng, hắt những vệt sáng lấp lánh lên từng góc phố sầm uất của thành phố Điện tử. Roboco vội vàng rảo bước chạy thật nhanh trên vỉa hè hướng về phía cửa hàng bách hóa tổng hợp, cô nàng vừa chạy vừa liên tục đưa mắt nhìn vào màn hình đồng hồ đeo tay đầy lo lắng: ``Trời đất ơi, trễ thế này rồi không biết trong tiệm người ta có còn sót lại hộp cà ri nào không nữa đây nhỉ? Hy vọng là vận may của mình hôm nay tốt, cửa hàng vẫn chưa bán hết sạch hàng\dots''

Ngay khi vừa bước chân vào bên trong không gian ấm áp, thơm phức của cửa hàng bách hóa, ánh mắt của nàng Người máy lập tức sáng rực lên đầy phấn khích khi phát hiện ra thứ mà bản thân hằng mong chờ suốt cả ngày hôm nay đang được đặt ngay ngắn, gọn gàng trên kệ hàng trưng bày: ``\textbf{Đây rồi! Thần may mắn mỉm cười với mình thật rồi này!}''

\img{4-3g}

Trước mặt cô chính là hộp cà ri thương hiệu ``Sina Dam (loại siêu cay) - hương vị thơm ngon đến mức bùng nổ vị giác! - phiên bản giới hạn''. Chẳng thèm suy nghĩ hay đắn đo lấy một giây, cô nàng lập tức vươn tay giật lấy hộp cà ri, ôm chặt nó vào lòng hệt như vừa mới tìm thấy một món báu vật vô giá của cuộc đời, rồi hăm hở bước nhanh về phía quầy thanh toán tiền.

Thế nhưng, ngay khi vừa mới tiến lại gần quầy thu ngân, một giọng đọc đơn điệu, đều đều nhưng lại vô cùng quen thuộc bỗng chốc vang lên đón khách: ``Dạ, xin kính chào quý khách. Em có thể giúp gì được cho chị không ạ?''

Roboco ngạc nhiên ngẩng đầu lên nhìn thì vô cùng kinh ngạc khi nhận ra người đang đứng sau quầy thu ngân chính là Pekora. Cô lúc này đang diện bộ đồng phục nhân viên cửa hàng, đôi mắt vàng cam dường như đã hoàn toàn mất đi sự sống và sự linh hoạt thường ngày, gương mặt lộ rõ vẻ mệt mỏi, chán chường tột độ của một người đang phải làm ca làm việc đêm muộn, tay lờ đờ cầm chiếc máy quét tính tiền.

``Ồ, Pekora-chan! Hóa ra em làm việc ở đây à?''

\img{4-3h}

Pekora khẽ thở dài một tiếng đầy ngao ngán, uể oải đáp ngắn gọn bằng chất giọng cộc lốc: ``Em làm ca tối thôi, peko.''

Nàng Người máy mỉm cười rạng rỡ, cúi thấp người xuống đặt hộp cà ri báu vật lên mặt quầy kính: ``À, ra vậy. Thế tính tiền nhanh cho chị đi\~{}.''

Pekora nhanh nhẹn đưa máy quét qua phần mã vạch, nhưng vẻ mặt vẫn giữ nguyên nét lạnh nhạt, mệt mỏi đặc trưng của nhân viên ca đêm. Sau khi máy quét xong và màn hình hiển thị giá cả hiện lên, cô nàng báo giá bằng giọng đều đều cộc lốc: ``Tổng cộng của chị hết\dots''

Roboco nhanh chóng rút tiền ra thanh toán hóa đơn, không quên nở một nụ cười tươi tắn cảm ơn cô em đồng nghiệp dễ thương: ``Chị cảm ơn em nhiều nhé, Pekora-chan.''

Pekora hời hợt gật đầu đáp lễ cho xong chuyện: ``Chúc chị có bữa tối ngon miệng, peko.'' Nói rồi, cô nàng lại lờ đờ hướng mắt về phía trước, tiếp tục cam chịu quay trở lại với công việc thu ngân vô cùng buồn tẻ của mình.

Cầm hộp cà ri trong tay, Roboco vừa rảo bước đi bộ quay trở về văn phòng vừa không kìm nén được nụ cười đầy phấn khích trên môi. Cô nàng Người máy lẩm bẩm đầy hào hứng trong lòng: ``\textit{Tuyệt vời quá, cuối cùng thì sau bao nhiêu sóng gió và khó nhọc, mình cũng đã có thể chính thức được thưởng thức món cà ri này rồi! Lại còn săn được hẳn bản giới hạn đặc biệt nữa chứ, quả này ngon quên sầu luôn!}''

\ct

Kuon đang ngồi lướt điện thoại ở bàn họp lớn khẽ ngẩng đầu lên nhìn khi nghe thấy tiếng bước chân vội vã: ``À, em về rồi đấy à Roboco-chan. Nhanh thế.''

Thế nhưng, lúc này Roboco dường như hoàn toàn chẳng mảy may bận tâm đến lời chào hỏi hay sự tồn tại của cậu Tổng. Với gương mặt ngập tràn vẻ hân hoan và sướng khoái tột độ, cô nhanh chóng chạy ùa tới chiếc bàn họp lớn, đặt hộp cà ri báu vật xuống mặt bàn rồi vội vàng dùng tay bóc vỏ, mở nắp hộp ra. Đôi mắt cô nàng sáng rực lên đầy háo hức hệt như một đứa trẻ chuẩn bị được khám phá một rương kho báu thực thụ: ``Cuối cùng thì\dots''

Để rồi, ngay khi nắp hộp vừa được bật mở ra, toàn bộ nụ cười rạng rỡ cùng vẻ mặt ngập tràn hạnh phúc của nàng Người máy bỗng chốc hoàn toàn biến mất sạch sẽ, thay vào đó là sự bối rối xen lẫn sững sờ tột độ. Đập vào mắt cô bên trong hộp cà ri hoàn toàn không có nước sốt hay thịt bò nào cả, mà là\dots\ hai người lão niên nhỏ bé cổ quái, chiều cao chưa đầy nửa đốt ngón tay, đang đứng lù lù bên trong.

Hai người ấy trông giống hệt như các nhân vật bước ra từ một câu chuyện cổ tích thần thoại Nhật Bản, diện bộ trang phục kimono giản dị và ánh mắt hiền từ. Điều kỳ lạ là hai cụ già tí hon đang đứng thản nhiên trò chuyện trò chuyện rôm rả với nhau, hoàn toàn ngó lơ sự tồn tại của một cô nàng Người máy khổng lồ đang há hốc mồm đứng thẫn thờ nhìn họ đầy khó hiểu kia.

\img{4-3i}

Họ lẩm bẩm những câu nói đầy trầm hùng hệt như đang cùng nhau đọc lên một lời sấm truyền tiên tri cổ xưa:

Lão ông khẽ vuốt chòm râu bạc trắng của mình, trầm ngâm nói: ``Ngày hôm đó, loài người nhỏ bé đã được điềm báo trước\dots''

Lão bà đứng bên cạnh cũng khẽ gật gù phụ họa tiếp lời đầy huyền bí: ``\dots về nỗi sợ hãi khi một thế lực khổng lồ cai quản họ.''

Nàng Người máy chớp mắt, không tin vào những gì mình đang nhìn thấy. Cô lắp bắp hỏi đầy hoang mang: ``\dots Hả? Họ đang tự tiện nói cái gì bên trong cái hộp ăn của tôi thế này cơ chứ\dots''

\img{4-3j}

Kuon ngồi đối diện chứng kiến toàn bộ cảnh tượng bóc hộp nãy giờ liền lập tức bật cười lớn nghiêng ngả, tay ôm chặt lấy bụng cười đến mức không thể kiềm chế được: ``Ôi trời đất ơi Roboco-chan ơi! Đúng thật là Robo-PON chính hiệu mà! Em lại bị người ta lừa mua phải hàng giả, hàng rởm nữa rồi con bé ngốc này ơi! Làm gì có cái thương hiệu cà ri chính hãng nào trên đời lại mang cái tên kỳ dị là 'cà ri Sina Dam' được cơ chứ hả em!?''

Nghe những lời châm chọc giễu cợt đầy phũ phàng của cậu Tổng, Roboco giận tím mặt. Cô nàng dứt khoát dùng lực đóng sầm nắp hộp cà ri lại một cách vô cùng dứt khoát, nhất quyết không thèm hé môi thốt thêm bất kỳ lời nào nữa. Cô lặng lẽ đặt chiếc hộp rởm lên mặt bàn, đôi mắt đờ đẫn, vô hồn: ``Thôi thì, cà ri chắc để sau vậy\dots''

Tiếng cười trêu chọc của cậu vẫn chưa chịu dừng lại, khiến cô nàng càng thêm phần ngượng ngùng và bối rối tột độ: ``Cái tên thương hiệu nghe đã thấy sặc mùi lừa đảo rồi, thế mà em vẫn cắm đầu vào mua cho bằng được, để rồi bây giờ khui hộp ra lại gặp ngay hai ông cụ tí hon đứng đọc tiên tri thế kia nữa chứ. Hóa ra đây chính là món 'cà ri đặc biệt' mà em hằng mong đợi suốt cả ngày hôm nay đấy à Roboco-chan?''

\img{4-3k}

Năm đĩa cà ri nóng hổi được xếp ngay ngắn trên bàn, tỏa hương nghi ngút kích thích mọi giác quan. Cả năm người họ cùng nhau ngồi xuống, chắp tay thành kính trước ngực và đồng thanh hô vang đầy hào hứng: ``Mời cả nhà ăn cơm\~{}!''

Ngay khi dứt lời, họ đồng loạt cầm thìa lên để thưởng thức bữa tối muộn sau một ngày dài đầy hỗn loạn. Không ngoài dự đoán, món cà ri tại gia này có hương vị đậm đà và vừa miệng đến kinh ngạc, nhưng người tỏ ra hạnh phúc và thỏa mãn nhất chắc chắn không ai khác ngoài Roboco. Kể từ trưa đến giờ, cô nàng Người máy chỉ biết ``hít khí trời'' và nhịn đói rã rời để chạy đôn chạy đáo làm nhiệm vụ giải cứu hết củ khoai tây lại đến thùng hành tây, nên lúc này mỗi thìa cà ri nóng hổi đưa vào miệng đều mang lại cảm giác ngọt ngào vô hạn.

Fubuki vừa ăn vừa vui vẻ múc thêm phần nước sốt đậm đà trên đĩa của mình, gật gù tấm tắc khen ngợi không tiếc lời: ``Trời đất ơi, cà ri ngon xuất sắc luôn ấy! Sự ngọt dịu từ đống hành tây được xào kỹ đúng là đỉnh cao mà. Cảm ơn Roboco-senpai và Kuon-senpai nhiều nha!''

Kuon khẽ dùng thìa xắn một miếng thịt bò mềm, ăn một cách vô cùng từ tốn và điềm đạm, mỉm cười đồng tình: ``Quả nhiên là tự tay vào bếp chuẩn bị nguyên liệu rồi nấu nướng ở nhà thì lúc nào ăn cũng thấy ngon miệng và yên tâm hơn hẳn so với đống hộp ăn liền trôi nổi ngoài bách hóa.''

Matsuri ở bên cạnh cũng không chịu thua kém, đôi tay máy của cô nàng liên tục hoạt động, xúc cơm lia lịa rồi hào hứng chen ngang: ``Chuẩn luôn anh ơi! Mà em nghĩ nồi cà ri này ngon đến vậy là vì bên trong có chứa cả 'tình yêu' cùng công sức đóng góp nguyên liệu của tất cả mọi người tụi mình đấy nhé\~{}!''

Trái ngược hoàn toàn với bầu không khí vui tươi, ăn uống hăng say của bốn người còn lại, Pekora ngồi ở góc bàn ăn một cách vô cùng chậm rãi và thầm lặng. Đôi tai thỏ dài trên đầu rũ cụp xuống, cô nàng vừa nhai vừa rơm rớm nước mắt đầy thương cảm nhìn chằm chằm vào những lát cà rốt cam rực trên đĩa cà ri của mình, thút thít lẩm bẩm đầy oán trách: ``Ăn thì ngon thật đấy peko\dots\ nhưng mà nghĩ lại thấy bất công ghê á. Rõ ràng là Roboco-senpai không nên tự tiện giật phắt phụ kiện cà rốt quý giá trên tóc của em như thế chứ, peko\dots''

Chứng kiến vẻ mặt mếu máo tội nghiệp nhưng miệng vẫn nhai nhồm nhoàm không ngừng của cô em thỏ, nàng Rô-bốt ngượng ngùng gãi gãi má, nở nụ cười đầy ái ngại chữa thẹn: ``Chị xin lỗi mà, Pekora-chan\~{}. Tại lúc đó tình thế cấp bách quá, trong bếp lại chẳng còn mẩu cà rốt nào cả. Để ngày mai chị chạy ra tiệm mua đền cho em hẳn một rổ cà rốt siêu to, siêu tươi khác bù lại nhé?''

Cậu Tổng quản lý đứng bên cạnh nghe vậy liền bật cười lớn đầy sảng khoái, vỗ vỗ vai cô nàng Người máy: ``Coi như đây là bài học nhớ đời cho em đấy nhé Roboco-chan, đừng bao giờ dại dột mà đi trộm đồ của nàng thỏ đanh đá này. Nhưng mà thôi, dù có bị mắng một chút nhưng đổi lại được một đĩa cà ri ngon lành thế này thì cũng hoàn toàn xứng đáng đấy chứ nhỉ?''

Roboco nghe cậu Tổng trêu chọc chỉ biết cúi đầu gãi gãi tóc, cười khì khì đầy ngốc nghếch: ``Chắc là thế nhỉ\dots?''

Tất cả mọi người đều bật cười với phản ứng rối bời của cô nàng ngốc nghếch đó. Buổi tối hôm ấy, cả năm người họ đã cùng nhau thưởng thức một bữa tiệc cà ri vô cùng ấm cúng, no nê và sạch sẽ, và quan trọng nhất là một món ăn hoàn toàn chất lượng, không hề chứa bất kỳ yếu tố lừa gạt hay xuất hiện cụ già tí hon nào đọc sấm truyền tiên tri cả.

\ct

Sau khi dọn dẹp sạch sẽ chén đĩa, cả nhóm kéo nhau ra ngồi vây quanh chiếc bàn họp lớn quen thuộc. Lúc này, mỗi người đều đang nâng niu trên tay một cốc trà xanh nóng hổi, bốc khói nghi ngút do đích thân Kuon pha chế. Hương trà thơm nhẹ thoang thoảng lan tỏa trong không gian tĩnh lặng của buổi đêm, mang lại một cảm giác nhẹ nhàng, thư thái và bình yên đến lạ kỳ sau một ngày dài hoạt động mệt mỏi. Tiếng cười nói nhỏ nhẹ, vui vẻ thỉnh thoảng lại vang lên, lấp đầy mọi ngóc ngách của căn phòng họp nhỏ.

Fubuki khẽ nhấp một ngụm trà ấm, ánh mắt mơ màng nhìn ra phía cửa sổ lớn hắt ánh đèn đường neon rực rỡ, cảm thán: ``Nghĩ lại thì hôm nay đúng là một ngày đầy ắp những bất ngờ thú vị nhỉ. Ai mà ngờ được chỉ từ một cơn thèm ăn đột xuất và một hộp cà ri Sina Dam rởm lại có thể dẫn tụi mình đến một bữa ăn chung vui vẻ và ấm áp đến thế này cơ chứ.''

``Quá chuẩn luôn chị ơi!'' Matsuri hào hứng hưởng ứng ngay lập tức, hai mắt sáng lên đầy nhiệt huyết: ``Lần sau tụi mình lại cùng nhau tụ tập bày trò nấu món gì khác đi! Em đề xuất lần tới tụi mình làm một nồi lẩu sukiyaki siêu to khổng lồ đi, ăn vào mùa xuân này thì còn gì bằng nữa!''

Pekora ngồi bên cạnh vẫn đang ôm cốc trà ấm, tai thỏ khẽ giật giật, cô nàng khẽ liếc xéo Roboco một cái rõ sắc nhưng khóe môi thì đã không kìm được mà hơi cong lên, nhỏ giọng tuyên bố đầy kiêu kỳ: ``Muốn ăn sukiyaki hay gì cũng được peko, nhưng nhớ là cấm tuyệt đối không được ai tự tiện cướp đồ ăn hay phụ kiện trên đầu của em nữa đâu đấy nha, peko! Tạm thời tha lỗi cho chị Roboco lần này thôi đấy!''

Roboco nghe lời cảnh báo đanh thép ấy liền chống cằm lên mặt bàn, nở một nụ cười ngại ngùng nhưng tràn đầy hạnh phúc, nhẹ nhàng thủ thỉ: ``Rồi mà, chị hứa từ nay sẽ không nghịch ngợm thế nữa đâu. Nhưng mà\dots\ mọi người thấy không? Dù cho ngày hôm nay tụi mình có gặp phải bao nhiêu rắc rối dở khóc dở cười đi chăng nữa, thì cuối cùng chúng ta vẫn kết thúc bằng những nụ cười vui vẻ thế này. Chẳng phải điều đó mới là quan trọng nhất sao?''

Lời chia sẻ đầy chân thành từ một cô nàng thường ngày vốn vô cùng ngốc nghếch và hay bày trò phá phách bỗng chốc khiến cả phòng họp chìm vào một khoảng lặng bình yên. Cả Kuon, Fubuki, Matsuri lẫn Pekora đều khẽ nhìn nhau rồi mỉm cười gật đầu, thầm đồng cảm sâu sắc với suy nghĩ ấy.

Kuon nâng cốc trà lên nhấp thêm một ngụm nhỏ, ánh mắt đầy ấm áp nhìn quanh gương mặt của những cô gái đang tràn đầy năng lượng tươi trẻ trước mặt, khẽ lên tiếng đầy tự hào: ``Hiếm khi nào anh thấy Roboco-chan nói được một câu triết lý và đúng đắn đến thế này đấy nhé. Những khoảnh khắc bình dị và vui vẻ thế này chính là lý do khiến cho công ty thần tượng của chúng ta luôn trở nên đặc biệt. Đâu chỉ là nơi làm việc áp lực, đây còn là mái nhà chung để mọi người cùng nhau kết nối, sẻ chia và tìm thấy niềm vui thực sự sau những giờ phút bận rộn.''

Ánh đèn vàng dịu nhẹ, ấm áp từ trần nhà chiếu xuống, bao trùm lên toàn bộ không gian phòng họp, làm cho cả căn phòng như được nhuộm một màu vàng nhạt đầy yên ả, thanh bình. Roboco nhìn quanh những gương mặt thân thuộc, cảm nhận được hơi ấm lan tỏa từ cốc trà truyền vào lòng bàn tay, rồi khẽ thốt lên bằng giọng điệu vô cùng nhẹ nhàng, thì thầm hệt như một lời tự sự đầy xúc động từ tận đáy lòng: ``Có lẽ\dots\ chính bầu không khí gia đình ấm áp này mới là điều khiến em cảm thấy no lòng và hạnh phúc nhất, ấm áp hơn cả việc ăn được một hộp cà ri nóng hổi rất nhiều\dots''

Cả căn phòng bỗng chốc lại rộn vang lên những tiếng cười đùa giòn giã đầy hạnh phúc của các cô gái, kéo dài mãi không dứt. Một buổi tối bình dị, đơn sơ nhưng ngập tràn những tiếng cười vui vẻ và những kỷ niệm đáng nhớ đã chính thức khép lại một ngày dài đầy rẫy những điều bất ngờ thú vị của văn phòng \hll.
 hăm hở chạy thẳng vào trong gian bếp nhỏ của văn phòng mới, mở toang các ngăn tủ để tìm kiếm bột gia vị cà ri cùng những dụng cụ nấu nướng cần thiết. Phần cơm trắng thì đã có sẵn gạo thơm lưu kho của công ty lo liệu, thế nhưng còn về phần rau củ ăn kèm\dots\ cô nàng bỗng chợt nảy ra một ý tưởng vô cùng táo bạo khi nghĩ ngay tới những người đồng nghiệp thân thiết của mình.

Cô nàng nhẩm tính trong đầu danh sách các nguyên liệu cần thu thập: ``\textit{Khoai tây thì lát nữa mình qua phòng Matsuri-chan xin một ít. Hành tây thì Fubuki-chan ban nãy vẫn còn cất giữ cả một thùng siêu to kia kìa. Còn cà rốt ư\dots\ À! Trên đầu của Pekora-chan lúc nào chẳng giắt sẵn hai củ cà rốt to đùng làm phụ kiện cơ chứ, qua mượn tạm em ấy là xong chuyện ngay thôi mà!}''

Nghĩ là làm, cô nàng lập tức tràn đầy năng lượng chạy ra ngoài để thực hiện kế hoạch đi thu thập nguyên liệu. Nhìn bóng lưng hăng hái hoạt bát một cách kinh ngạc của cô nàng Người máy, Kuon chỉ biết khẽ lắc đầu cười thầm đầy chiều chuộng: ``\textit{Haiz, đúng là sức mạnh của một tâm hồn ăn uống có khác, có thể khiến một cô nàng lười biếng bỗng chốc trở nên năng nổ đến mức này cơ đấy\dots\ Nhưng mà tự tay vào bếp nấu ăn thế này cũng là một trải nghiệm vô cùng xứng đáng và ý nghĩa đấy chứ.}''

\ct

Matsuri sau khi nghe thấy Roboco muốn xin khoai tây liền vô cùng vui vẻ mang đến một túi khoai tây tròn trịa, tươi ngon (tất nhiên đây hoàn toàn là khoai tây thật dùng để nấu ăn chứ không phải là những chiếc máy báo thức ngụy trang), cô nàng hớn hở cười nói: ``Đây nè Roboco-senpai ơi! Đống khoai tây này là phần thưởng em giành được từ một gian hàng trò chơi ở hội chợ tuần trước đấy ạ. Em tặng lại hết cho chị luôn nè!''

Nhìn túi khoai tây ú cụ, béo múp mà cô nàng Lễ hội nhiệt tình trao tặng, Roboco vô cùng hào hứng reo lên đầy cảm kích: ``Chị cảm ơn em nhiều nhiều lắm nhé Matsuri-chan! Có bấy nhiêu đây khoai tây là đã quá đủ để làm một nồi cà ri thịnh soạn cho cả nhà rồi!''

Ngay trước khi nàng Người máy quay người rời khỏi phòng để tiếp tục đi tìm nguyên liệu khác, Matsuri còn không quên tinh nghịch nháy mắt một cái trêu chọc đầy ẩn ý: ``Nhưng nhớ là sau khi nấu xong xuôi thì chị phải đặc biệt mời em một đĩa cà ri siêu to đấy nha\~{}!''

Đáp lại lời đề nghị dễ thương của cô em, Roboco vui vẻ giơ ngón tay cái hướng lên trời tạo dáng đồng ý đầy dứt khoát.

\ct

Roboco sau đó tìm thấy Fubuki đang ngồi cặm cụi soát lại đống hành tây lưu kho ngoài nắng hôm trước ở góc hành lang. Cô nàng liền tiến lại gần, mỉm cười mở lời xin xỏ: ``Fubuki-chan, tình hình là chị đang rất cần vài củ hành tây để nấu nồi cà ri ăn tối. Em có thể cho chị xin mấy củ được không em?''

Nhìn thấy cô nàng Người máy hôm nay bỗng dưng lại có hứng thú vào bếp một cách vô cùng lạ thường, nàng Cáo trắng liền vui vẻ cúi xuống bới bới lục lọi bên trong chiếc thùng gỗ của mình: ``Dạ được chứ chị, để em kiểm tra xem trong đống này còn sót lại củ nào tươi ngon dùng được không đã nhé\dots''

Fubuki cẩn thận lật qua lật lại từng củ hành tây một cách vô cùng kỹ lưỡng dưới ánh đèn, vừa lọc vừa lẩm bẩm đầy chuyên nghiệp: ``Hừm, củ này trông héo quá không được\dots\ củ này thì bám nhiều đất bẩn quá cũng không ổn luôn\dots''

Sau một hồi loay hoay lục lọi tìm kiếm, cô nàng bỗng reo lên đầy hào hứng: ``À, đây rồi!'' Cô chìa ra ba củ hành tây tròn trĩnh, đưa cho Roboco.

``Chị lấy mấy củ này nhé. Nhưng nhớ là đừng để chúng ngoài nắng lâu quá, không là lại 'mọc râu' như lần trước đấy!'' nàng Cáo trắng nói với tông giọng châm chọc.

Nàng Rô-bốt nghe vậy liền ngại ngùng cười trừ đầy chữa thẹn: ``Chị biết rồi\~{} Mà, chị tưởng chị phải là người nói câu này với em nhỉ? Em là người phơi nó ra mà.''

\ct

Pekora lúc này vừa mới chính thức hoàn thành xong ca làm việc đêm mệt mỏi ở bách hóa bèn lê lết từng bước chân rã rời đầy mệt mỏi quay trở về văn phòng mới: ``Mệt quá, peko\dots\ Lại còn chưa có gì bỏ bụng nữa chứ\dots'' cô nàng than vãn, lết thân người uể oải về văn phòng.

Bất ngờ, cô bị Roboco nhảy ra chặn đứng đường đi. Gương mặt của cô nàng Người máy lúc này lại đang hiện rõ một sự hào hứng, phấn khích vô cùng rạng rỡ, hoàn toàn trái ngược với tâm trạng mệt mỏi hiện tại của cô nàng Thỏ kia.

Nàng Người máy nhìn chăm chăm vào đầu Pekora, mắt của cô nàng sáng rực lên đầy ranh mãnh: ``Pekora-chan, cà rốt trên đầu em thật đẹp\dots''

Vẫn với sự mệt mỏi và uể oải đó, nàng Thỏ nhìn cô bằng ánh mắt ngập trạng sự cảnh giác: ``\dots Peko? Chị đang nhìn cái gì trên đầu em vậy\dots?''

Nàng Người máy mỉm cười ngọt ngào, hai bàn tay máy của cô đang ngoe nguẩy tiến sát về phía mấy củ cà rốt đang được giắt vô cùng khéo léo bên trong hai lọn tóc tết: ``Pekora-chan, cho chị xin một củ nhé\~{}''

Nhìn thấy bàn tay máy đang có ý định xâm phạm phụ kiện tóc yêu quý, Pekora lập tức nhảy lùi lại phía sau phòng thủ, cố gắng giữ lấy những củ cà rốt trên tóc của mình như thể đó là vật bất ly thân: ``K-Không đời nào có chuyện đó đâu nha, peko! Đây là vật dụng cá nhân của em đấy! Cấm chị tuyệt đối không được phép chạm vào, peko!''

Không kịp để Pekora dứt câu phản đối, nàng Rô-bốt nhanh tay lẹ mắt rút phắt ngay vài củ cà rốt từ tóc của nàng Thỏ khiến đối phương hoàn toàn không kịp trở tay ứng phó. Pekora giật nảy mình, la toáng lên đầy kinh hoàng: ``\textbf{Gyaaaah! Chị đang tự tiện làm cái trò cướp bóc trắng trợn gì thế này chị Roboco-senpai, peko!?}''

Roboco không thèm trả lời cô em, chỉ ôm chặt lấy chiến lợi phẩm rồi quay đầu chạy biến về phía nhà bếp với tốc độ cực nhanh, còn không quên ngoảnh đầu lại vẫy tay chào: ``Cảm ơn em nhé, Pekora-chan\~{}!''

Nàng Thỏ liền ngã lăn ra sàn nhà phòng họp, không ngừng giãy giụa lăn lộn ăn vạ như đứa trẻ bị giật mất đồ chơi: ``\textbf{Trả lại cà rốt cho em đi mà, peko!! Đó là cà rốt của em mà\dots}'' Cô nàng vừa khóc lóc thảm thiết vừa la toáng lên để gây sự chú ý và cầu xin sự công bằng từ phía Kuon và Fubuki đang đứng gần đó.

Nhìn cô nàng Thỏ xanh đáng thương đang gào mồm lên khóc lóc, hai tay ôm lấy phần lọn tóc giờ đã hoàn toàn trống trơn, cậu Tổng quản lý Kuon khẽ bật cười thành tiếng đầy hoài niệm, tự nhẩm: ``\textit{Giống cái hồi mà mình với Koro-chan và Noel-chan trêu em ấy hồi Comiket mùa đông (xem lại ÉPISODE 34, phần Truyện phụ, quyển số Ba) nhỉ.}''

Trong khi đó nàng Cáo trắng Fubuki ngỡ ngác và bối rối trước hành động như một đứa trẻ nít bị giành mất đồ chơi của nàng Thỏ, vội vàng chạy lại đỡ cô nàng dậy hỏi han: ``P-Pekora-chan!? Sao thế?''

``Roboco-senpai\dots\ lấy mất cà rốt của em rồi\dots'' nàng Thỏ xanh nói trong tiếng nấc nghẹn ngào, rồi tiếp tục nằm lăn ra sàn ăn vạ tiếp.

Cuối cùng, cả Cáo trắng lẫn cậu Tổng quản lý phải mất đến hơn mười phút dỗ dành thì nàng Thỏ mới chịu nín khóc và thôi không tiếp tục nằm ăn vạ nữa.

\ct

Sau khoảng bốn mươi phút cùng nhau hì hục xoay xở trong bếp, bầu không khí mát mẻ của buổi tối muộn dường như đã bị xua tan hoàn toàn bởi mùi hương thơm nức, béo ngậy của bột gia vị cà ri. Nồi nước sốt vàng óng, sánh mịn đang sôi lục bục trên bếp, tỏa ra hơi ấm dễ chịu và mang theo hương thơm ngào ngạt của những miếng thịt bò hầm mềm, quyện lẫn vị ngọt tự nhiên của hành tây Fubuki, những lát khoai tây tròn trịa của Matsuri, và cả những khoanh cà rốt tươi rói vừa được tịch thu từ đầu Pekora.

Đứng trước thành quả tự tay mình chuẩn bị, Roboco không giấu nổi niềm hào hứng khôn xiết, đôi mắt sáng rỡ hệt như đứa trẻ vừa nhận được món quà mơ ước. Cô nàng Người máy nhanh nhẹn múc cơm nóng ra đĩa, khéo léo chan nước sốt cà ri bò bốc khói nghi ngút lên trên rồi cẩn thận bày biện năm suất ăn vô cùng thịnh soạn, đẹp mắt ra bàn ăn lớn.

\img{4-3k}

Năm đĩa cà ri nóng hổi được xếp ngay ngắn trên bàn, tỏa hương nghi ngút kích thích mọi giác quan. Cả năm người họ cùng nhau ngồi xuống, chắp tay thành kính trước ngực và đồng thanh hô vang đầy hào hứng: ``Mời cả nhà ăn cơm\~{}!''

Ngay khi dứt lời, họ đồng loạt cầm thìa lên để thưởng thức bữa tối muộn sau một ngày dài đầy hỗn loạn. Không ngoài dự đoán, món cà ri tại gia này có hương vị đậm đà và vừa miệng đến kinh ngạc, nhưng người tỏ ra hạnh phúc và thỏa mãn nhất chắc chắn không ai khác ngoài Roboco. Kể từ trưa đến giờ, cô nàng Người máy chỉ biết ``hít khí trời'' và nhịn đói rã rời để chạy đôn chạy đáo làm nhiệm vụ giải cứu hết củ khoai tây lại đến thùng hành tây, nên lúc này mỗi thìa cà ri nóng hổi đưa vào miệng đều mang lại cảm giác ngọt ngào vô hạn.

Fubuki vừa ăn vừa vui vẻ múc thêm phần nước sốt đậm đà trên đĩa của mình, gật gù tấm tắc khen ngợi không tiếc lời: ``Trời đất ơi, cà ri ngon xuất sắc luôn ấy! Sự ngọt dịu từ đống hành tây được xào kỹ đúng là đỉnh cao mà. Cảm ơn Roboco-senpai và Kuon-senpai nhiều nha!''

Kuon khẽ dùng thìa xắn một miếng thịt bò mềm, ăn một cách vô cùng từ tốn và điềm đạm, mỉm cười đồng tình: ``Quả nhiên là tự tay vào bếp chuẩn bị nguyên liệu rồi nấu nướng ở nhà thì lúc nào ăn cũng thấy ngon miệng và yên tâm hơn hẳn so với đống hộp ăn liền trôi nổi ngoài bách hóa.''

Matsuri ở bên cạnh cũng không chịu thua kém, đôi tay máy của cô nàng liên tục hoạt động, xúc cơm lia lịa rồi hào hứng chen ngang: ``Chuẩn luôn! Mà em nghĩ nồi cà ri này ngon đến vậy là vì bên trong có chứa cả `tình yêu' cùng công sức đóng góp nguyên liệu của tất cả mọi người tụi mình đấy nhé\~{}!''

Trái ngược hoàn toàn với bầu không khí vui tươi, ăn uống hăng say của bốn người còn lại, Pekora ngồi ở góc bàn ăn một cách vô cùng chậm rãi và thầm lặng. Đôi tai thỏ dài trên đầu rũ cụp xuống, cô nàng vừa nhai vừa rơm rớm nước mắt đầy thương cảm nhìn chằm chằm vào những lát cà rốt cam rực trên đĩa cà ri của mình, thút thít lẩm bẩm đầy oán trách: ``Ăn thì ngon thật đấy peko\dots\ nhưng mà nghĩ lại thấy bất công ghê á. Rõ ràng là Roboco-senpai không nên tự tiện giật phắt phụ kiện cà rốt quý giá trên tóc của em như thế chứ, peko\dots''

Chứng kiến vẻ mặt mếu máo tội nghiệp nhưng miệng vẫn nhai nhồm nhoàm không ngừng của cô em thỏ, nàng Rô-bốt ngượng ngùng gãi gãi má, nở nụ cười đầy ái ngại chữa thẹn: ``Chị xin lỗi mà, Pekora-chan\~{}. Tại lúc đó tình thế cấp bách quá, trong bếp lại chẳng còn mẩu cà rốt nào cả. Để ngày mai chị chạy ra tiệm mua đền cho em hẳn một rổ cà rốt siêu to, siêu tươi khác bù lại nhé?''

Cậu Tổng quản đứng bên cạnh nghe vậy liền bật cười lớn đầy sảng khoái, vỗ vỗ vai cô nàng Người máy: ``Coi như đây là bài học nhớ đời cho em đấy nhé Roboco-chan, đừng bao giờ dại dột mà đi trộm đồ của nàng thỏ đanh đá này. Nhưng mà thôi, dù có bị mắng một chút nhưng đổi lại được một đĩa cà ri ngon lành thế này thì cũng hoàn toàn xứng đáng đấy chứ nhỉ?''

Roboco nghe cậu trêu chọc chỉ biết cúi đầu gãi gãi tóc, cười khì khì đầy ngốc nghếch: ``Chắc là thế nhỉ\dots?''

Tất cả mọi người đều bật cười với phản ứng rối bời của cô nàng ngốc nghếch đó. Buổi tối hôm ấy, cả năm người họ đã cùng nhau thưởng thức một bữa tiệc cà ri vô cùng ấm cúng, no nê và sạch sẽ, và quan trọng nhất là một món ăn hoàn toàn chất lượng, không hề chứa bất kỳ yếu tố lừa gạt hay xuất hiện cụ già tí hon nào đọc sấm truyền tiên tri cả.

\ct

Sau khi dọn dẹp sạch sẽ chén đĩa, cả nhóm kéo nhau ra ngồi vây quanh chiếc bàn họp lớn quen thuộc. Lúc này, mỗi người đều đang nâng niu trên tay một cốc trà xanh nóng hổi, bốc khói nghi ngút do đích thân Kuon pha chế. Hương trà thơm nhẹ thoang thoảng lan tỏa trong không gian tĩnh lặng của buổi đêm, mang lại một cảm giác nhẹ nhàng, thư thái và bình yên đến lạ kỳ sau một ngày dài hoạt động mệt mỏi. Tiếng cười nói nhỏ nhẹ, vui vẻ thỉnh thoảng lại vang lên, lấp đầy mọi ngóc ngách của căn phòng họp nhỏ.

Fubuki khẽ nhấp một ngụm trà ấm, ánh mắt mơ màng nhìn ra phía cửa sổ lớn hắt ánh đèn đường neon rực rỡ, cảm thán: ``Nghĩ lại thì hôm nay đúng là một ngày đầy ắp những bất ngờ thú vị nhỉ. Ai mà ngờ được chỉ từ một cơn thèm ăn đột xuất và một hộp cà ri Sina Dam rởm lại có thể dẫn tụi mình đến một bữa ăn chung vui vẻ và ấm áp đến thế này cơ chứ.''

``Quá chuẩn luôn chị ơi!'' Matsuri hào hứng hưởng ứng ngay lập tức, hai mắt sáng lên đầy nhiệt huyết: ``Lần sau tụi mình lại cùng nhau tụ tập bày trò nấu món gì khác đi! Em đề xuất lần tới tụi mình làm một nồi lẩu sukiyaki siêu to khổng lồ đi, ăn vào mùa xuân này thì còn gì bằng nữa!''

Pekora ngồi bên cạnh vẫn đang ôm cốc trà ấm, tai thỏ khẽ giật giật, cô nàng khẽ liếc xéo Roboco một cái rõ sắc nhưng khóe môi thì đã không kìm được mà hơi cong lên, nhỏ giọng tuyên bố đầy kiêu kỳ: ``Muốn ăn sukiyaki hay gì cũng được peko, nhưng nhớ là cấm tuyệt đối không được ai tự tiện cướp đồ ăn hay phụ kiện trên đầu của em nữa đâu đấy nha, peko! Tạm thời tha lỗi cho chị boco lần này thôi đấy!''

Roboco nghe lời cảnh báo đanh thép ấy liền chống cằm lên mặt bàn, nở một nụ cười ngại ngùng nhưng tràn đầy hạnh phúc, nhẹ nhàng thủ thỉ: ``Rồi mà, chị hứa từ nay sẽ không nghịch ngợm thế nữa đâu. Nhưng mà\dots\ mọi người thấy không? Dù cho ngày hôm nay tụi mình có gặp phải bao nhiêu rắc rối dở khóc dở cười đi chăng nữa, thì cuối cùng chúng ta vẫn kết thúc bằng những nụ cười vui vẻ thế này. Chẳng phải điều đó mới là quan trọng nhất sao?''

Lời chia sẻ đầy chân thành từ một cô nàng thường ngày vốn vô cùng ngốc nghếch và hay bày trò phá phách bỗng chốc khiến cả phòng họp chìm vào một khoảng lặng bình yên. Cả Kuon, Fubuki, Matsuri lẫn Pekora đều khẽ nhìn nhau rồi mỉm cười gật đầu, thầm đồng cảm sâu sắc với suy nghĩ ấy.

Kuon nâng cốc trà lên nhấp thêm một ngụm nhỏ, ánh mắt đầy ấm áp nhìn quanh gương mặt của những cô gái đang tràn đầy năng lượng tươi trẻ trước mặt, khẽ lên tiếng đầy tự hào: ``Hiếm khi nào anh thấy Roboco-chan nói được một câu triết lý và đúng đắn đến thế này đấy nhé. Những khoảnh khắc bình dị và vui vẻ thế này chính là lý do khiến cho công ty thần tượng của chúng ta luôn trở nên đặc biệt. Đâu chỉ là nơi làm việc áp lực, đây còn là mái nhà chung để mọi người cùng nhau kết nối, sẻ chia và tìm thấy niềm vui thực sự sau những giờ phút bận rộn.''

Ánh đèn vàng dịu nhẹ, ấm áp từ trần nhà chiếu xuống, bao trùm lên toàn bộ không gian phòng họp, làm cho cả căn phòng như được nhuộm một màu vàng nhạt đầy yên ả, thanh bình. Roboco nhìn quanh những gương mặt thân thuộc, cảm nhận được hơi ấm lan tỏa từ cốc trà truyền vào lòng bàn tay, rồi khẽ thốt lên bằng giọng điệu vô cùng nhẹ nhàng, thì thầm hệt như một lời tự sự đầy xúc động từ tận đáy lòng: ``Có lẽ\dots\ chính bầu không khí gia đình ấm áp này mới là điều khiến em cảm thấy no lòng và hạnh phúc nhất, ấm áp hơn cả việc ăn được một hộp cà ri nóng hổi rất nhiều\dots''

Cả căn phòng bỗng chốc lại rộn vang lên những tiếng cười đùa giòn giã đầy hạnh phúc của các cô gái, kéo dài mãi không dứt. Một buổi tối bình dị, đơn sơ nhưng ngập tràn những tiếng cười vui vẻ và những kỷ niệm đáng nhớ đã chính thức khép lại một ngày dài đầy rẫy những điều bất ngờ thú vị của văn phòng \hll.

\end{document}